\chapter{Conclusion générale}

\section{Bilan du projet}
Ce stage de deux mois au sein de l'entreprise Bewize a été une expérience 
particulièrement enrichissante, tant sur le plan technique que professionnel. 
Il m'a permis de participer activement à un projet concret et de contribuer 
de manière significative au développement d'une fonctionnalité réelle 
intégrée dans l'écosystème technique de Bewize.

L'objectif principal de mon stage était de développer une fonctionnalité 
complète de recueil de feedbacks et de réclamations pour l'application 
mobile Bewize. Cet objectif a été pleinement atteint : la solution que 
j'ai développée permet désormais aux utilisateurs de l'application mobile 
de transmettre facilement leurs retours, qu'il s'agisse de feedbacks 
positifs ou de réclamations, directement depuis l'application mobile 
via une page WebView intuitive et accessible. Ces informations sont 
ensuite centralisées dans le dashboard monitor interne de Bewize, où 
l'équipe peut les consulter, les analyser et les traiter de manière 
organisée et efficace.

\section{Récapitulatif des réalisations}
Au cours de ce stage, j'ai accompli plusieurs réalisations concrètes 
et significatives, couvrant l'ensemble des composants de la solution :

\begin{itemize}
    \item \textbf{Formulaire WebView React JS :} j'ai développé une 
    interface web légère, responsive et intuitive permettant aux 
    utilisateurs de soumettre leurs feedbacks et réclamations en 
    quelques clics, avec une validation complète des données saisies 
    et un retour visuel clair après l'envoi.
    
    \item \textbf{Intégration mobile React Native :} j'ai réussi 
    à intégrer le formulaire WebView dans l'application mobile Bewize 
    existante, avec l'ajout d'un bouton d'accès dédié et une navigation 
    fluide entre les différents écrans de l'application.
    
    \item \textbf{Back-end Spring Boot :} j'ai développé un ensemble 
    d'APIs REST robustes et bien documentées, gérant la réception, 
    la validation, le traitement et la persistance des données de 
    feedbacks et réclamations dans la base de données PostgreSQL.
    
    \item \textbf{Base de données PostgreSQL :} j'ai conçu et mis 
    en place un modèle de données adapté pour stocker les feedbacks 
    et réclamations de manière fiable, avec les contraintes d'intégrité 
    nécessaires.
    
    \item \textbf{Dashboard monitor React JS :} j'ai développé une 
    interface web interne complète permettant à l'équipe Bewize de 
    consulter, filtrer et traiter les retours reçus, avec un système 
    de gestion des statuts et un tableau de bord statistiques.
    
    \item \textbf{Conteneurisation Docker :} j'ai mis en place 
    l'environnement Docker avec docker-compose pour garantir la 
    cohérence entre les environnements de développement et de 
    production et simplifier le déploiement.
    
    \item \textbf{Tests complets :} j'ai réalisé une stratégie 
    de test complète couvrant les tests unitaires, les tests 
    d'intégration et les tests fonctionnels bout en bout, avec 
    un taux de réussite de 100\% sur l'ensemble des 39 tests réalisés.
\end{itemize}

\section{Compétences acquises}

\subsection{Compétences techniques}
Ce stage m'a permis d'acquérir et de consolider de nombreuses 
compétences techniques directement applicables dans le domaine 
du développement logiciel professionnel :

\begin{itemize}
    \item \textbf{Développement React JS et React Native :} j'ai 
    acquis une maîtrise approfondie du développement d'interfaces 
    web et mobiles avec React JS et React Native, notamment la 
    gestion des composants, des états, des appels API et de la navigation.
    
    \item \textbf{Développement back-end Spring Boot :} j'ai compris 
    et pratiqué le développement d'APIs REST avec Spring Boot, incluant 
    la gestion des entités JPA, des repositories, des services et 
    des contrôleurs.
    
    \item \textbf{Bases de données PostgreSQL :} j'ai pratiqué la 
    conception de modèles de données relationnels, la configuration 
    de PostgreSQL et son administration via pgAdmin.
    
    \item \textbf{Conteneurisation Docker :} j'ai découvert et 
    pratiqué la conteneurisation avec Docker et docker-compose 
    pour l'orchestration de services multiples.
    
    \item \textbf{Versioning Git et GitHub :} j'ai consolidé mes 
    pratiques de versioning avec Git, notamment la gestion des 
    branches selon la méthodologie Gitflow et la collaboration 
    via les pull requests et les code reviews.
    
    \item \textbf{Tests logiciels :} j'ai découvert et pratiqué 
    les différents niveaux de tests (unitaires, intégration, 
    fonctionnels) avec JUnit, Mockito, Jest et Postman.
\end{itemize}

\subsection{Compétences méthodologiques}
Au-delà des compétences purement techniques, ce stage m'a permis 
d'acquérir des compétences méthodologiques précieuses :

\begin{itemize}
    \item \textbf{Méthodologie Agile Scrum :} j'ai participé 
    activement aux rituels Scrum (Daily Scrum, Sprint Planning, 
    Sprint Review, Retrospective) et j'ai développé une compréhension 
    approfondie du framework Scrum dans un contexte professionnel réel.
    
    \item \textbf{Travail en équipe :} j'ai vécu l'expérience 
    de la collaboration avec une équipe pluridisciplinaire composée 
    d'un Product Manager, d'un Tech Lead et de développeurs, dans 
    un environnement professionnel exigeant et stimulant.
    
    \item \textbf{Gestion du temps et des priorités :} j'ai appris 
    à gérer efficacement mon temps dans un contexte de stage à durée 
    limitée, avec des objectifs précis et des délais à respecter.
    
    \item \textbf{Communication professionnelle :} j'ai développé 
    mes capacités de communication technique avec l'équipe de 
    développement et de communication fonctionnelle avec le Product 
    Manager lors des présentations et démonstrations.
\end{itemize}

\subsection{Découverte du marketing digital}
Conformément au planning de mon stage, la dernière phase a été 
consacrée à la découverte du marketing digital chez Bewize. Cette 
expérience complémentaire m'a permis de comprendre le lien essentiel 
entre le développement produit et la stratégie marketing d'une 
entreprise numérique. Cette vision globale du fonctionnement d'une 
startup comme Bewize constitue un atout précieux pour la suite de 
mon parcours professionnel.

\section{Apport pour Bewize}
La fonctionnalité que j'ai développée représente une valeur ajoutée 
concrète et mesurable pour Bewize et ses utilisateurs. Elle comble 
un manque identifié dans l'écosystème existant en offrant :

\begin{itemize}
    \item \textbf{Une meilleure expérience utilisateur :} les 
    utilisateurs de l'application mobile disposent désormais d'un 
    canal simple, rapide et accessible pour exprimer leurs feedbacks 
    et réclamations, renforçant leur sentiment d'être écoutés et 
    pris en compte par Bewize.
    
    \item \textbf{Une centralisation des retours :} l'équipe Bewize 
    dispose désormais d'une vue centralisée et organisée de tous 
    les retours utilisateurs dans le dashboard monitor, facilitant 
    leur analyse et leur traitement.
    
    \item \textbf{Une amélioration continue du produit :} les 
    feedbacks et réclamations collectés constituent une source 
    précieuse d'informations pour orienter les futurs développements 
    et améliorer continuellement l'application mobile Bewize.
    
    \item \textbf{Une meilleure réactivité :} le système de gestion 
    des statuts du dashboard monitor permet à l'équipe de suivre 
    précisément l'état de traitement de chaque réclamation et de 
    répondre plus rapidement aux utilisateurs insatisfaits.
\end{itemize}

\section{Perspectives et améliorations futures}
Bien que la solution que j'ai développée réponde pleinement aux 
besoins exprimés dans le cahier des charges, plusieurs pistes 
d'amélioration et d'évolution ont été identifiées pour les 
versions futures :

\begin{itemize}
    \item \textbf{Système de notifications :} mettre en place un 
    système de notifications push pour alerter l'équipe Bewize 
    dès qu'une nouvelle réclamation est reçue, permettant une 
    réactivité encore plus grande.
    
    \item \textbf{Réponse aux utilisateurs :} développer une 
    fonctionnalité permettant à l'équipe Bewize de répondre 
    directement aux utilisateurs ayant soumis une réclamation, 
    fermant ainsi la boucle de communication.
    
    \item \textbf{Catégorisation automatique :} intégrer un 
    mécanisme de catégorisation automatique des retours par 
    thématique (problème technique, demande de fonctionnalité, 
    bug, etc.) pour faciliter leur traitement et leur analyse.
    
    \item \textbf{Export des données :} ajouter une fonctionnalité 
    d'export des feedbacks et réclamations au format Excel ou CSV 
    pour faciliter l'analyse statistique par l'équipe produit.
    
    \item \textbf{Tableau de bord avancé :} enrichir le tableau 
    de bord statistiques avec des graphiques d'évolution dans le 
    temps, des indicateurs de satisfaction et des comparaisons 
    entre périodes.
    
    \item \textbf{Internationalisation :} adapter le formulaire 
    WebView et le dashboard monitor pour supporter plusieurs langues, 
    en anticipation d'une expansion internationale de l'application 
    Bewize.
\end{itemize}

\section{Mot de fin}
Ce stage chez Bewize a constitué une expérience professionnelle 
complète et particulièrement formatrice. Il m'a permis de découvrir 
concrètement le fonctionnement d'une équipe produit dans un 
environnement agile, de participer à un projet de développement 
logiciel réel de bout en bout, et de développer des compétences 
techniques et méthodologiques directement applicables dans ma 
future carrière professionnelle.

La richesse de cette expérience réside non seulement dans les 
compétences techniques que j'ai acquises, mais aussi dans la 
compréhension globale du développement produit : comment identifier 
un besoin métier, le transformer en spécifications techniques, 
développer une solution, la tester rigoureusement et la livrer 
dans un environnement agile en collaboration avec une équipe 
pluridisciplinaire.

Bewize m'a offert un cadre de stage idéal, alliant rigueur technique, 
bienveillance de l'équipe et projets stimulants. Cette expérience 
constitue une base solide pour la suite de mon parcours professionnel 
dans le domaine du développement logiciel.